\begin{corollary}[Proved under stated hypotheses: elliptic omitted singularities]
\label{D:elliptic-surface-boundary}
Let $S$ be a normal proper connected surface over an algebraically
closed field, and let $Y\subset S$ be finite.  Assume that every
singular point in $Y$ admits a resolution whose reduced exceptional
locus is one smooth elliptic curve.  Then
$\pi_1(S\setminus Y)$ is topologically finitely generated.
The singularities retained in $S\setminus Y$ are unrestricted.
For a possibly nonnormal proper scheme of dimension at most two,
the same conclusion holds if the finite inverse image of $Y$ on
each normal surface component satisfies this condition and $Y$
has codimension at least two in every component.
\end{corollary}
\begin{proof}
Choose a resolution $q:X\to S$ which agrees with the specified
single-curve resolutions around the singular points of $Y$, is an
isomorphism at its regular points, and resolves the remaining
singularities arbitrarily.  These choices glue: take affine
neighborhoods of each of the finitely many singular points
containing no other singular point; their resolutions are all
isomorphisms on the overlaps with the regular locus.
The resulting $X$ is smooth and proper over $k$, hence projective
by \cite[Lemma 54.16.11, tag 0C5P]{StacksDSurfaceResolution}.
Put $U=S\setminus Y$ and $U'=q^{-1}(U)$.

A connected finite \'etale cover of the normal $U$ stays connected
after the proper birational pullback to $U'$.  Its normalization
over $S$ is finite and \'etale away from $Y$.  At every divisorial
component of $X\setminus U'$ the local situation is the elliptic
case of Proposition~\ref{D:single-curve-residue-tameness}; all these
extensions are therefore tame.  The regular omitted points remain
codimension two and give no divisorial boundary conditions.
Consequently the surjection $\pi_1(U')\to\pi_1(U)$ factors
through a surjection
\[
              Q(X,X\setminus U')\longrightarrow\pi_1(U).
\]
The source is finitely generated by
Proposition~\ref{D:fixed-boundary-positive}, proving the claim.

In the nonnormal case we pass to the finite disjoint union of
normalizations of the reduced irreducible components.  The proved
surface case applies to each corresponding open; any components
of dimension at most one have no omitted point and are proper.
Their fundamental groups are finitely generated.  The finite
proper descent argument of Proposition~\ref{D:fixed-boundary-positive}
then gives finite generation downstairs.
\end{proof}
